\documentclass[varwidth]{article}
\usepackage[a4paper, margin=2.5cm]{geometry}

\usepackage{amsmath}
\usepackage{amssymb}

\begin{document}

{ \Large Aufgabe 1 - LB. S. 14 Nr. 1.}

geg.: $f(x) = -\frac1x ;~ g(x) = 2x-3 ;~ h(x) = \sqrt{x+3}-3$

\vspace*{5mm}

a)

$
\begin{array}{llll}

    ~    & x = -2 & x = 0,1 & x = 78 \\
    f(x) & 0,5 & -10 & \approx -0,013 \\
    g(x) & -7 & -2,8 & 153 \\
    h(x) & -2 & \approx -1,239 & 6 \\

\end{array}
$

\vspace*{5mm}

b)

$ D_f: x\in\mathbb R, x\neq0 $

$ D_g: x\in\mathbb R $

$ D_h: x\in\mathbb R, x\ge-3 $

\vspace*{5mm}

c)

geg.: $ P(1|-1);~ Q(5,5|8) $

\vspace*{5mm}

$
f(x_P) = f(1) = -\frac11 \\~\\
\boxed{ -\frac11 = -1 = y_P }
$

\vspace*{8mm}

$
g(x_P) = g(1) = 2\cdot1-3 \\~\\
\boxed{ 2\cdot1-3 = -1 = y_P }
$

\vspace*{8mm}

$
h(x_P) = h(1) = \sqrt{1+3}-3 \\~\\
\boxed{ \sqrt{1+3}-3 = -1 = y_P }
$

\vspace*{12mm}

$
f(x_Q) = f(5,5) = -\frac1{5,5} \\~\\
\boxed{ -\frac1{5,5} = -\frac{2}{11} \neq y_Q }
$

\vspace*{8mm}

$
g(x_Q) = g(5,5) = 2\cdot5,5-3 \\~\\
\boxed{ 2\cdot5,5-3 = 8 = y_Q }
$

\vspace*{8mm}

$
h(x_Q) = h(5,5) = \sqrt{5,5+3}-3 \\~\\
\boxed{ \sqrt{8,5}-3 \approx -0,085 \neq y_Q }
$

\vspace*{20mm}

\pagebreak
{ \Large Aufgabe 2 - Terme und Funktionen: Grundaufgaben}

{ \large 1. Bestimmen Sie die Nullstellen folgender Funktionen. }

\vspace*{5mm}

a)

$
\begin{array}{l}
f(x) = (x+3)^2 - 1 \\
f(x) = x^2 + 6x + 8 \\~\\
x_0^2 + 6x_0 + 8 = 0 \\
\rightarrow \boxed{x_{01,~ 02} = -3 \pm 1}
\end{array}
$

\vspace*{5mm}

b)

$
\begin{array}{l}
f(x) = 2x-5 \\~\\
2x_0-5 = 0 \\
\rightarrow \boxed{x_0 = 2,5}
\end{array}
$

\vspace*{5mm}

c) 

$
\begin{array}{l}
f(x) = (x+1)(3x-12) \\
f(x) = 3x^2 - 9x - 12 \\~\\
3x_0^2 - 9x_0 - 12 = 0 \\
x_0^2 - 3x_0 - 4 = 0 \\
\rightarrow \boxed{ x_{01,~ 02} = 1,5 \pm 2,5 }
\end{array}
$

\vspace*{5mm}

d)

$
\begin{array}{l}
f(x) = x(x+1)(x-3) \\
f(x) = x^3 - 2x^2 - 3x \\~\\
x_0^3 - 2x_0^2 - 3x_0 = 0 \\
\rightarrow \boxed{ x_{01}=-1;~ x_{02}=0;~ x_{03}=3 }
\end{array}
$

\vspace*{5mm}

e)

$
\begin{array}{l}
f(x) = x^3 - 4x^2 + 3x \\~\\
x_0^3 - 4x_0^2 + 3x_0 = 0 \\
\rightarrow \boxed{ x_{01} = 0;~ x_{02} = 1;~ x_{03} = 3 }
\end{array}
$

\vspace*{5mm}

{\large 2. Klammern Sie aus.}

\vspace*{5mm}

a) $\displaystyle 10x^3 - 6x^2 + 8x = 2x \cdot ( 5x^2 - 3x + 4 ) $

b) $\displaystyle 6u^4r^3 + 9u^2r - 3ur^2 = 3ur\cdot( 2u^3r^2 + 3u - r ) $

c) $\displaystyle 3b^2 + 2b + 5 = b^2 \cdot ( 3 + \frac2b + \frac5{b^2} ) $

\vspace*{5mm}

\pagebreak
{\large 3. Bestimmen Sie die Funktionsgleichung der Geraden.}

\vspace*{5mm}

a)

$
\begin{array}{l}
    m = 0,5; A(-5|-2) \\
    f(x) = mx+n \\~\\
    f(x_A) = y_A \\~\\
    0,5\cdot-5 + n = -2 \\
    \rightarrow n = 0,5 \\~\\
    \Rightarrow \boxed{f(x) = 0,5x + 0,5}
\end{array}
$

\vspace*{5mm}

b)

$
\begin{array}{l}
A(0|3); B(2|1) \\
f(x) = mx+n \\~\\
f(x_A) = y_A \\
f(x_B) = y_B \\~\\
\rightarrow
\begin{array}{rl}
    \text{I.} & 0m + n = 3 \\
    \text{II.} & 2m + n = 1
\end{array}
\\~\\
0m + n = 3 \rightarrow n = 3
\\~\\
2m + n = 1 \\
2m + 3 = 1 \rightarrow m = -1
\\~\\
\Rightarrow \boxed{f(x) = -1x + 3}
\end{array}
$

\vspace*{5mm}

{\large 4. Bestimmen Sie die Funktionsgleichung der Exponentialfunktion der Form $f(x)=a^x$ durch den Punkt}

a) $P(1|5)$ ~~~~~ b) $R(-1|4)$

\vspace*{5mm}

a)

$
\begin{array}{l}
f(x) = a^x \\
a^1 = 5 \rightarrow a = 5 \\~\\
\Rightarrow \boxed{f(x) = 5^x}
\end{array}
$

\vspace*{5mm}

b)

$
\begin{array}{l}
f(x) = a^x \\
a^{-1} = 4 \rightarrow a = \frac14 \\~\\
\Rightarrow \boxed{f(x) = \frac14^x}
\end{array}
$

\vspace*{5mm}

\pagebreak
{\large 5. Berechnen Sie die Schnittpunkte der folgenden Funktionen.}

\vspace*{5mm}

a)

$
\begin{array}{l}
f(x) = x+1 \\
g(x) = 3x-5
\\~\\
x+1 = 3x-5 \\
\rightarrow x = 3
\\~\\
f(3) = 4 \\
\Rightarrow \boxed{S(3|4)}
\end{array}
$

\vspace*{5mm}

b)

$
\begin{array}{l}
f(x) = x^2 + 2 \\
g(x) = 3x^2-4x-14
\\~\\
x^2 + 2 = 3x^2 - 4x - 14 ~~|~ -x^2-2\\
2x^2-4x-16 = 0
\\
\rightarrow x_{1,2} = 1\pm3
\\~\\
f(x_1) = 18 \\
f(x_2) = 6 \\
\Rightarrow \boxed{ S_1(4|18); S_2(-2|6) }
\end{array}
$

\vspace*{20mm}
{ \Large Aufgabe 3 - LB. S. 192 Nr. 1.}

\vspace*{5mm}

geg.:

\makebox[3mm]{}
\vbox{
    Funktionen $f_a; f_b; \dots; f_e; f_f$ der Form $\displaystyle f(x) = \frac{u(x)}{v(x)}$ \par
    Funktion $\text{grad}(f)$, definiert als der Grad des Parameterterms (Bsp.: $\text{grad}(x^3+2x^2+5) = 3$)
}

\vspace*{5mm}

\begin{tabular}{|l|l|l|l|}
    \hline
    Funktion $f$ & $\displaystyle\lim_{x\to\infty}f(x)$ & $\displaystyle\lim_{x\to-\infty}f(x)$ & Asymptote \\
    \hline
    $\displaystyle f_a(x) = \frac{x^2-5x}{3x^2+x}$
        & $\text{grad}(u_a(x)) = \text{grad}(v_a(x)) \rightarrow \frac13$
        & $\text{grad}(u_a(x)) = \text{grad}(v_a(x)) \rightarrow \frac13$
        & $y = \frac13$ \\
    % \hline
    &&&\\
    $\displaystyle f_b(x) = \frac{4x^2+22}{1+3x^3}$
        & $\text{grad}(u_b(x)) < \text{grad}(v_b(x)) \rightarrow 0$
        & $\text{grad}(u_b(x)) < \text{grad}(v_b(x)) \rightarrow 0$
        & $y = 0$ \\
    % \hline
    &&&\\
    $\displaystyle f_c(x) = \frac{28x^3-x}{7x^3+x^2}$
        & $\text{grad}(u_c(x)) = \text{grad}(v_c(x)) \rightarrow 4$
        & $\text{grad}(u_c(x)) = \text{grad}(v_c(x)) \rightarrow 4$
        & $y = 4$ \\
    % \hline
    &&&\\
    $\displaystyle f_d(x) = \frac{8x^3-5x^2}{x^2-2x^3}$
        & $\text{grad}(u_d(x)) = \text{grad}(v_d(x)) \rightarrow -4$
        & $\text{grad}(u_d(x)) = \text{grad}(v_d(x)) \rightarrow -4$
        & $y = -4$ \\
    % \hline
    &&&\\
    $\displaystyle f_e(x) = \frac{x^2+0,2x^4}{1+x^3}$
        & $\text{grad}(u_e(x)) = \text{grad}(v_e(x))+1 \rightarrow \infty$
        & $\text{grad}(u_e(x)) = \text{grad}(v_e(x))+1 \rightarrow -\infty$
        & $y = \frac x5$ \\
    \hline
\end{tabular}

\end{document}